% GENERATED FILE. Edit canonical lessons, then run npm run generate:books.
% canonical-source-hash: fnv1a64:b3c507ab657078c2
% canonical-lessons: GU-C390-pak, GU-C390-paripakva, GU-C390-ujjad, GU-C390-ugadvu, GU-C390-upajvu

\chapter{Making It Grow}
\label{ch:gu-making-it-grow}
\hlchaptermodality{390}

\begin{chapteropening}
\textbf{By the end of this chapter:} \emph{I can say what will grow if we plant and wait.}

Complete the last lesson of chapter 390: I can say what will grow if we plant and wait.
\end{chapteropening}

\section[\texorpdfstring{\gu{પાક} (pāk)}{પાક (pāk)}]{\gu{પાક} (pāk) — a crop}

\label{lesson:GU-C390-pak}



\begin{quote}
Before the new one: say the Gujarati for to hesitate, then the Gujarati for to waver.
\end{quote}

\subsection*{You'll want to know: \gu{પાક}}
\textbf{\gu{પાક}} — \emph{pāk} — \textquotedblleft{}a crop\textquotedblright{}.

\begin{quote}
\gu{જો} \gu{સમયસર} \gu{વરસાદ} \gu{પડશે} \gu{તો} \gu{આ} \gu{વરસે} \gu{પાક} \gu{સારો} \gu{થશે}.
\end{quote}

\begin{quote}
\emph{jo samaysar varsād paḍśe to ā varse pāk sāro thaśe.}
\end{quote}

\begin{quote}
If the rain comes on time, the crop will be good this year.
\end{quote}

A masculine noun.

\begin{cousinweb}
\textbf{\gu{પાક}} (\emph{pāk}) shares its source with \textbf{\gu{પાકવું}} (\emph{pākvũ}), \textquotedblleft{}to ripen\textquotedblright{}: CDIAL traces both to Sanskrit \emph{pakvá}, \textquotedblleft{}cooked, ripe\textquotedblright{}, which in the Atharvaveda already meant \textquotedblleft{}ripe corn\textquotedblright{} (CDIAL 7621). CDIAL glosses the Gujarati noun only as \textquotedblleft{}sugar boiled for conserves\textquotedblright{}, so \textquotedblleft{}a crop\textquotedblright{} is a sense it does not give.
\end{cousinweb}

\begin{grammarlens}[title={What you've built}]
One more word for saying what happens if.
\end{grammarlens}

\subsection*{Your turn}
\begin{itemize}
\raggedright
  \item \emph{Say it:} \emph{pāk}
  \item \emph{Say it:} \emph{pāk}, once more
  \item \emph{Say it:} \emph{jo samaysar varsād paḍśe to ā varse pāk sāro thaśe.} — If the rain comes on time, the crop will be good this year.
  \item \emph{Recall:} say the Gujarati for to hesitate, then the Gujarati for to waver, then say \emph{pāk} again
\end{itemize}

\begin{tcolorbox}[breakable,colback=teal!4,colframe=teal!35!black,title={Before you move on}]
What does \gu{પાક} mean? (\textquotedblleft{}A crop\textquotedblright{}.) Say it once more.
\end{tcolorbox}
\section[\texorpdfstring{\gu{પરિપક્વ} (paripakva)}{પરિપક્વ (paripakva)}]{\gu{પરિપક્વ} (paripakva) — mature}

\label{lesson:GU-C390-paripakva}



\begin{quote}
Before the new one: say the Gujarati for to waver, then the Gujarati for a crop.
\end{quote}

\subsection*{You'll want to know: \gu{પરિપક્વ}}
\textbf{\gu{પરિપક્વ}} — \emph{paripakva} — \textquotedblleft{}mature\textquotedblright{}.

\begin{quote}
\gu{જો} \gu{તે} \gu{થોડો} \gu{વધારે} \gu{પરિપક્વ} \gu{હોત} \gu{તો} \gu{આવી} \gu{ભૂલ} \gu{ન} \gu{કરત}.
\end{quote}

\begin{quote}
\emph{jo te thoḍo vadhāre paripakva hot to āvī bhūl na karat.}
\end{quote}

\begin{quote}
If he were a little more mature, he would not make such a mistake.
\end{quote}

\begin{cousinweb}
\textbf{\gu{પરિપક્વ}} (\emph{paripakva}) is a tatsama, borrowed from Sanskrit \emph{paripakva}, \textquotedblleft{}completely cooked, fully ripe\textquotedblright{} (Monier-Williams).
\end{cousinweb}

\begin{grammarlens}[title={What you've built}]
One more word for saying what happens if.
\end{grammarlens}

\subsection*{Your turn}
\begin{itemize}
\raggedright
  \item \emph{Say it:} \emph{paripakva}
  \item \emph{Say it:} \emph{paripakva}, once more
  \item \emph{Say it:} \emph{jo te thoḍo vadhāre paripakva hot to āvī bhūl na karat.} — If he were a little more mature, he would not make such a mistake.
  \item \emph{Recall:} say the Gujarati for to waver, then the Gujarati for a crop, then say \emph{paripakva} again
\end{itemize}

\begin{tcolorbox}[breakable,colback=teal!4,colframe=teal!35!black,title={Before you move on}]
What does \gu{પરિપક્વ} mean? (\textquotedblleft{}Mature\textquotedblright{}.) Say it once more.
\end{tcolorbox}
\section[\texorpdfstring{\gu{ઉજ્જડ} (ujjaḍ)}{ઉજ્જડ (ujjaḍ)}]{\gu{ઉજ્જડ} (ujjaḍ) — barren}

\label{lesson:GU-C390-ujjad}



\begin{quote}
Before the new one: say the Gujarati for a crop, then the Gujarati for mature.
\end{quote}

\subsection*{You'll want to know: \gu{ઉજ્જડ}}
\textbf{\gu{ઉજ્જડ}} — \emph{ujjaḍ} — \textquotedblleft{}barren\textquotedblright{}.

\begin{quote}
\gu{જો} \gu{આપણે} \gu{જંગલ} \gu{કાપીશું} \gu{તો} \gu{આ} \gu{ધરતી} \gu{ઉજ્જડ} \gu{થઈ} \gu{જશે}.
\end{quote}

\begin{quote}
\emph{jo āpaṇe jaṅgal kāpīśũ to ā dhartī ujjaḍ thaī jaśe.}
\end{quote}

\begin{quote}
If we cut down the forest, this land will turn barren.
\end{quote}

\begin{cousinweb}
\textbf{\gu{ઉજ્જડ}} (\emph{ujjaḍ}) CDIAL lists the Gujarati word, \textquotedblleft{}waterless, barren\textquotedblright{}, under a reconstructed \emph{ujjaṭa-}, \textquotedblleft{}uprooted, laid waste\textquotedblright{} (CDIAL 1660b).
\end{cousinweb}

\begin{grammarlens}[title={What you've built}]
One more word for saying what happens if.
\end{grammarlens}

\subsection*{Your turn}
\begin{itemize}
\raggedright
  \item \emph{Say it:} \emph{ujjaḍ}
  \item \emph{Say it:} \emph{ujjaḍ}, once more
  \item \emph{Say it:} \emph{jo āpaṇe jaṅgal kāpīśũ to ā dhartī ujjaḍ thaī jaśe.} — If we cut down the forest, this land will turn barren.
  \item \emph{Recall:} say the Gujarati for a crop, then the Gujarati for mature, then say \emph{ujjaḍ} again
\end{itemize}

\begin{tcolorbox}[breakable,colback=teal!4,colframe=teal!35!black,title={Before you move on}]
What does \gu{ઉજ્જડ} mean? (\textquotedblleft{}Barren\textquotedblright{}.) Say it once more.
\end{tcolorbox}
\section[\texorpdfstring{\gu{ઉગાડવું} (ugāḍvũ)}{ઉગાડવું (ugāḍvũ)}]{\gu{ઉગાડવું} (ugāḍvũ) — to grow plants}

\label{lesson:GU-C390-ugadvu}



\begin{quote}
Before the new one: say the Gujarati for mature, then the Gujarati for barren.
\end{quote}

\subsection*{You'll want to know: \gu{ઉગાડવું}}
\textbf{\gu{ઉગાડવું}} — \emph{ugāḍvũ} — \textquotedblleft{}to grow plants\textquotedblright{}.

\begin{quote}
\gu{જો} \gu{તમે} \gu{આંગણામાં} \gu{ટામેટાં} \gu{ઉગાડશો} \gu{તો} \gu{બજારમાંથી} \gu{ખરીદવાં} \gu{નહીં} \gu{પડે}.
\end{quote}

\begin{quote}
\emph{jo tame ā̃gṇāmā̃ ṭāmeṭā̃ ugāḍśo to bajārmā̃thī kharīdvā̃ nahī̃ paḍe.}
\end{quote}

\begin{quote}
If you grow tomatoes in the courtyard, you will not have to buy them at the market.
\end{quote}

Past, with \gu{તેણે} (he or she …): \gu{તેણે} \gu{બગીચામાં} \gu{ફૂલ} \gu{ઉગાડ્યાં}, \emph{teṇe bagīcāmā̃ phūl ugāḍyā̃}.

\begin{cousinweb}
\textbf{\gu{ઉગાડવું}} (\emph{ugāḍvũ}) is the causative of \textbf{\gu{ઊગવું}} (\emph{ūgvũ}), \textquotedblleft{}to rise, to sprout\textquotedblright{}, from a reconstructed \emph{udgāti}, \textquotedblleft{}rises, comes forth\textquotedblright{} (CDIAL 1954).
\end{cousinweb}

\begin{grammarlens}[title={What you've built}]
One more word for saying what happens if.
\end{grammarlens}

\subsection*{Your turn}
\begin{itemize}
\raggedright
  \item \emph{Say it:} \emph{ugāḍvũ}
  \item \emph{Say it:} \emph{ugāḍvũ}, once more
  \item \emph{Say it:} \emph{jo tame ā̃gṇāmā̃ ṭāmeṭā̃ ugāḍśo to bajārmā̃thī kharīdvā̃ nahī̃ paḍe.} — If you grow tomatoes in the courtyard, you will not have to buy them at the market.
  \item \emph{Recall:} say the Gujarati for mature, then the Gujarati for barren, then say \emph{ugāḍvũ} again
\end{itemize}

\begin{tcolorbox}[breakable,colback=teal!4,colframe=teal!35!black,title={Before you move on}]
What does \gu{ઉગાડવું} mean? (\textquotedblleft{}To grow plants\textquotedblright{}.) Say it once more.
\end{tcolorbox}
\section[\texorpdfstring{\gu{ઊપજવું} (ūpajvũ)}{ઊપજવું (ūpajvũ)}]{\gu{ઊપજવું} (ūpajvũ) — to be produced (as a crop)}

\label{lesson:GU-C390-upajvu}



\begin{quote}
Before the new one: say the Gujarati for barren, then the Gujarati for to grow plants.
\end{quote}

\subsection*{You'll want to know: \gu{ઊપજવું}}
\textbf{\gu{ઊપજવું}} — \emph{ūpajvũ} — \textquotedblleft{}to be produced (as a crop)\textquotedblright{}.

\begin{quote}
\gu{જો} \gu{ખેતરમાં} \gu{પૂરતું} \gu{પાણી} \gu{હશે} \gu{તો} \gu{ઘણા} \gu{ઘઉં} \gu{ઊપજશે}.
\end{quote}

\begin{quote}
\emph{jo khetarmā̃ pūrtũ pāṇī haśe to ghaṇā ghaũ ūpajśe.}
\end{quote}

\begin{quote}
If the field has enough water, it will yield a lot of wheat.
\end{quote}

Past (it …): \gu{તે} \gu{ખેતરમાં} \gu{બાજરી} \gu{ઊપજી}, \emph{te khetarmā̃ bājrī ūpjī}.

\begin{cousinweb}
\textbf{\gu{ઊપજવું}} (\emph{ūpajvũ}) is inherited from Sanskrit \emph{útpadyatē}, \textquotedblleft{}arises, originates\textquotedblright{}; CDIAL lists the Gujarati verb, \textquotedblleft{}is produced\textquotedblright{} (CDIAL 1814).
\end{cousinweb}

\begin{grammarlens}[title={What you've built}]
One more word for saying what happens if.
\end{grammarlens}

\subsection*{Your turn}
\begin{itemize}
\raggedright
  \item \emph{Say it:} \emph{ūpajvũ}
  \item \emph{Say it:} \emph{ūpajvũ}, once more
  \item \emph{Say it:} \emph{jo khetarmā̃ pūrtũ pāṇī haśe to ghaṇā ghaũ ūpajśe.} — If the field has enough water, it will yield a lot of wheat.
  \item \emph{Recall:} say the Gujarati for barren, then the Gujarati for to grow plants, then say \emph{ūpajvũ} again
\end{itemize}

\begin{tcolorbox}[breakable,colback=teal!4,colframe=teal!35!black,title={Before you move on}]
What does \gu{ઊપજવું} mean? (\textquotedblleft{}To be produced (as a crop)\textquotedblright{}.) Say it once more.
\end{tcolorbox}
